% Zone „Das kennst du schon“ – aus bank/zuordnungen/zone.jsonl (zusammenbau v0.1)
\einheitenkopf[zone][Kennst du schon]{Das kennst du schon}
\setcounter{aufgabe}{0}
%% TODO Grundvorstellungs-Aufgabe als erste Hauptnummer der Zone (2.8) fehlt in der Bank

%% TODO Titel: Ich-kann-Satz fehlt in der Bank (Platzhalter: Kettenname) – Nr. 1
%% TODO Anweisung: ein Satz über der ersten Teilaufgabe fehlt in der Bank – Nr. 1
\begin{aufgabe}{Punkte im Koordinatensystem eintragen und ablesen (erster Quadrant)}
\swa{Trage den Punkt A(4|3) ins Koordinatensystem ein.}{\begin{ksys}[xmin=0,xmax=6,ymin=0,ymax=6]\end{ksys}}
\swz{Lies ab: Welche Koordinaten hat der Punkt C?}{\begin{ksys}[xmin=0,xmax=6,ymin=0,ymax=6,ablesen]\punkt{2.5}{4}{C}\end{ksys}}
\end{aufgabe}

%% TODO Titel: Ich-kann-Satz fehlt in der Bank (Platzhalter: Kettenname) – Nr. 2
%% TODO Anweisung: ein Satz über der ersten Teilaufgabe fehlt in der Bank – Nr. 2
\begin{aufgabe}{Multiplizieren und Dividieren mit Dezimalzahlen (Preis geteilt durch Stückzahl, Preis mal Stückzahl)}
%% TODO Darstellung für schwach fehlt in der Bank (zuordnungen-zone-f2-v1; Abschnitt „Für schwache Schüler“)
\swz{3 Hefte kosten je 1,50 €. Was kosten sie zusammen?}{}
%% TODO Darstellung für schwach fehlt in der Bank (zuordnungen-zone-f2-v3; Abschnitt „Für schwache Schüler“)
\swz{7 Kugeln Eis kosten je 1,20 €. Was kosten alle zusammen?}{}
\end{aufgabe}

%% TODO Titel: Ich-kann-Satz fehlt in der Bank (Platzhalter: Kettenname) – Nr. 3
%% TODO Anweisung: ein Satz über der ersten Teilaufgabe fehlt in der Bank – Nr. 3
\begin{aufgabe}{Vielfache und Teiler erkennen (zwölf ist das Dreifache von vier)}
%% TODO Darstellung für schwach fehlt in der Bank (zuordnungen-zone-f3-v1; Abschnitt „Für schwache Schüler“)
\swz{15 ist das Wievielfache von 5?}{}
%% TODO Darstellung für schwach fehlt in der Bank (zuordnungen-zone-f3-v3; Abschnitt „Für schwache Schüler“)
\swz{3 kg sind das Wievielfache von 0,5 kg?}{}
\end{aufgabe}

%% TODO Titel: Ich-kann-Satz fehlt in der Bank (Platzhalter: Kettenname) – Nr. 4
%% TODO Anweisung: ein Satz über der ersten Teilaufgabe fehlt in der Bank – Nr. 4
\begin{aufgabe}{Bruchteil einer Größe (die Hälfte, ein Viertel von 12 €)}
%% TODO Darstellung für schwach fehlt in der Bank (zuordnungen-zone-f4-v1; Abschnitt „Für schwache Schüler“)
\swz{Die Hälfte von 14 € – wie viel?}{}
%% TODO Darstellung für schwach fehlt in der Bank (zuordnungen-zone-f4-v3; Abschnitt „Für schwache Schüler“)
\swz{Ein Viertel von 3 € – wie viel?}{}
\end{aufgabe}

%% TODO Titel: Ich-kann-Satz fehlt in der Bank (Platzhalter: Kettenname) – Nr. 5
%% TODO Anweisung: ein Satz über der ersten Teilaufgabe fehlt in der Bank – Nr. 5
\begin{aufgabe}{Einheiten umrechnen (Cent ↔ Euro, Minuten ↔ Stunden, m ↔ km)}
%% TODO Darstellung für schwach fehlt in der Bank (zuordnungen-zone-f5-v1; Abschnitt „Für schwache Schüler“)
\swz{Wie viele Minuten sind 3 Stunden?}{}
%% TODO Darstellung für schwach fehlt in der Bank (zuordnungen-zone-f5-v3; Abschnitt „Für schwache Schüler“)
\swz{Wie viele Meter sind 4,2 km?}{}
\end{aufgabe}

%% TODO Titel: Ich-kann-Satz fehlt in der Bank (Platzhalter: Kettenname) – Nr. 6
%% TODO Anweisung: ein Satz über der ersten Teilaufgabe fehlt in der Bank – Nr. 6
\begin{aufgabe}{Vielfache und Teiler erkennen (zwölf ist das Dreifache von vier) – Fehler finden}
\swz[3]{Finn schreibt: „Von 7 € auf 28 € sind es 21 € mehr, also sind 28 € das 21-Fache von 7 €.“ Finde den Fehler und rechne richtig.}{}
\end{aufgabe}

%% TODO Titel: Ich-kann-Satz fehlt in der Bank (Platzhalter: Kettenname) – Nr. 7
%% TODO Anweisung: ein Satz über der ersten Teilaufgabe fehlt in der Bank – Nr. 7
\begin{aufgabe}{Vielfache und Teiler erkennen (zwölf ist das Dreifache von vier) – selbst rechnen}
%% TODO Darstellung für schwach fehlt in der Bank (zuordnungen-zone-f3-v6; Abschnitt „Für schwache Schüler“)
\swz{Ein Beet war 9 m² groß, jetzt ist es 45 m² groß. Das Wievielfache ist das?}{}
\end{aufgabe}
